\section{Conclusion}\label{conclusion}

Language models develop structured moral representations that go
beyond mere moral detection. By training independent linear probes
for each of six Moral Foundations Theory foundations and analyzing
the geometry of their weight vectors, we find that OLMo-2 1B and
OLMoE-1B-7B exhibit the \emph{integration} signature: foundation
directions are distinct (effective dimensionality \(= 5\), which rules
out collapse) yet share a common moral-salience component, shown by a
uniformly positive mean pairwise cosine (\(\approx 0.22\)--\(0.27\)) and a
leading principal component that captures \({\sim}0.38\) of the
directions' variance (vs.~\({\sim}0.18\) for random directions). The
positive cosine, not the effective dimensionality, is what
distinguishes integration from isolation. This
geometric structure is consistent across the two architectures
tested (dense vs.~MoE), emerges early in pre-training, and
stabilizes before probing accuracy saturates. However, the
inter-framework structure does not align with MFT's predicted
individualizing/binding grouping, indicating that the model's
moral taxonomy is empirically grounded in corpus statistics rather
than organized along human-theoretical lines.

Framework-specific fragility testing shows that the output dilution
effect is foundation-uniform: every foundation is more fragile in MoE
than in dense (a \({\sim}2.3\times\) per-foundation gap), with no
foundation reliably most or least robust within an architecture and no
significant binding/individualizing group difference.

Extending this geometric lens to moral dilemmas (scenarios
where two foundations conflict) shows partial compositionality:
dilemma representations share \({\sim}10\)\% of their variance with
the subspace spanned by component foundation directions (100\(\times\)
the null baseline), with near-balanced loading across both
conflicting foundations. The remaining \({\sim}90\)\% residual and the
complexity--fragility gradient (\(\sigma^*\): 4.72 \(\to\) 3.12 \(\to\)
2.90) indicate that dilemma representations encode conflict-specific
structure beyond their component foundations. This compositionality
pattern is preserved across architectures (dense and MoE).

The probe direction geometry methodology introduced here is
general. Any set of related concepts that can be isolated by
binary linear probes can be analyzed for inter-concept structure
via the same cosine similarity and clustering techniques. We
anticipate applications to other taxonomies of human values, to
political orientation, and to the internal organization of safety
training in instruction-tuned models.
